\section{The family of order \texorpdfstring{$4k+4$}{4k+4}}\label{sec:family4}
The second family in the article begins with graphs on eight, twelve and
sixteen vertices at strict levels one, two and three. All three occur as
nested induced subgraphs. From the sixteen-vertex seed onwards, the layer
construction adds four vertices at a time. We specify these graphs, justify
their finite reductions and obstructions, and then apply the general theorem.

\subsection{The sixteen-vertex seed and its profiles}\label{4k4:sec:seed}
We use three auxiliary types, so each layer adds four vertices. Let
\(V(A)=\{a_0,\ldots,a_{15}\}\), with the adjacency list in
Table~\ref{4k4:tab:Aadj}. It is symmetric, has no loops, and has 97 edges.
The profiles in Table~\ref{4k4:tab:profiles} have sizes 13, 14 and 15.
The type graph is
\begin{equation}\label{4k4:eq:R}
E(R)=\{\{0,1\},\{0,2\}\}.
\end{equation}
\begin{table}[H]
\centering\small
\caption{Complete adjacency list of the sixteen-vertex seed \(A\).}
\label{4k4:tab:Aadj}
\begin{tabular}{r l}\toprule
\(v\)&\(N_A(a_v)\), given by indices\\\midrule
0 & 2, 3, 5, 6, 7, 8, 9, 10, 11, 12, 13, 14\\
1 & 2, 3, 4, 6, 7, 8, 9, 10, 11, 12, 13, 14\\
2 & 0, 1, 3, 6, 7, 8, 9, 10, 11, 12, 13, 15\\
3 & 0, 1, 2, 4, 5, 8, 9, 10, 11, 12, 13, 15\\
4 & 1, 3, 5, 6, 7, 8, 10, 11, 12, 13, 14, 15\\
5 & 0, 3, 4, 6, 7, 9, 10, 11, 12, 13, 14, 15\\
6 & 0, 1, 2, 4, 5, 8, 9, 11, 12, 13, 14, 15\\
7 & 0, 1, 2, 4, 5, 8, 9, 10, 12, 13, 14, 15\\
8 & 0, 1, 2, 3, 4, 6, 7, 9, 10, 12, 14, 15\\
9 & 0, 1, 2, 3, 5, 6, 7, 8, 11, 13, 14, 15\\
10 & 0, 1, 2, 3, 4, 5, 7, 8, 11, 13, 14, 15\\
11 & 0, 1, 2, 3, 4, 5, 6, 9, 10, 12, 14, 15\\
12 & 0, 1, 2, 3, 4, 5, 6, 7, 8, 11, 14, 15\\
13 & 0, 1, 2, 3, 4, 5, 6, 7, 9, 10, 14, 15\\
14 & 0, 1, 4, 5, 6, 7, 8, 9, 10, 11, 12, 13, 15\\
15 & 2, 3, 4, 5, 6, 7, 8, 9, 10, 11, 12, 13, 14\\
\bottomrule\end{tabular}
\end{table}

\begin{table}[H]
\centering
\caption{Profiles specified by their omitted seed vertices. Every copy
of type \(j\) is adjacent to all remaining vertices of \(A\).}
\label{4k4:tab:profiles}
\begin{tabular}{r l r}\toprule
\(j\)&\(V(A)\setminus P_j\)&\(\card{P_j}\)\\\midrule
0 & \(\{a_8,a_{11},a_{14}\}\)&13\\
1 & \(\{a_{11},a_{15}\}\)&14\\
2 & \(\{a_{13}\}\)&15\\
\bottomrule\end{tabular}
\end{table}

The graph \(B=G_1\) adds \(p=c_0\) and \(x_0,x_1,x_2\) to \(A\).
The centre is adjacent to every vertex of \(A\) and to no auxiliary.
The neighbours of \(x_j\) in \(A\) are \(P_j\), and its auxiliary
adjacencies are specified by \(R\). Thus
\[
\card{V(B)}=20,\qquad \card{E(B)}=97+16+42+2=157.
\]
For the finite certificates, \(a_i\) has index \(i\), \(p=16\),
and \(x_j=17+j\). 
\subsection{The eight- and twelve-vertex prefixes}
The first two members are induced subgraphs of the same seed. Put
\[
\begin{aligned}
S_1&=\{2,4,5,8,9,10,12,13\},\\
S_2&=\{2,3,4,5,6,7,8,9,10,11,12,13\},
\end{aligned}
\]
and define
\[
H_1=A[\{a_i:i\in S_1\}],\qquad
H_2=A[\{a_i:i\in S_2\}],\qquad H_3=A.
\]
The direct reduction below establishes the strict level of the first prefix.
The recursive certificates then establish the strict level of the second
prefix and of the sixteen-vertex seed.
The two prefixes have 8 and 12 vertices, respectively.
The inclusions \(S_1\subset S_2\subset V(A)\) give induced inclusions.
We now verify the first prefix directly. Use numerical labels \(a_i=i\). Restricting the adjacency
table of \(A\) gives
\[
\begin{array}{c|l}
v&N_{H_1}(v)\\\hline
2&8,9,10,12,13\\
4&5,8,10,12,13\\
5&4,9,10,12,13\\
8&2,4,9,10,12\\
9&2,5,8,13\\
10&2,4,5,8,13\\
12&2,4,5,8\\
13&2,4,5,9,10
\end{array}
\]
\begin{proposition}\label{4k4:prop:manual-prefix}
The graph \(H_1\) belongs to \(\D_1\setminus\D_0\).
\end{proposition}
\begin{proof}
First consider \(H_1[N_{H_1}(5)]\), on \(\{4,9,10,12,13\}\).
The only neighbour of \(9\) there is \(13\), so \(9\to13\)
is a fold. The remaining graph on \(\{4,10,12,13\}\) is a cone
with apex \(4\). Consequently \(5\) can be deleted from \(H_1\)
at level one.

In \(H_1-5\), the neighbourhood of \(4\) is
\(\{8,10,12,13\}\). Within this neighbourhood, \(12\to8\)
is a fold, and the remaining path \(8,10,13\) is a cone with
apex \(10\). Thus \(4\) can next be deleted at level one.
The graph \(H_1-\{5,4\}\) is a cone with apex \(2\), proving
\(H_1\in\D_1\). The inner folds justify only the stated outer
deletion; they do not delete vertices from the outer graph.

To exclude level zero, Table~\ref{4k4:tab:manual-prefix-negative}
lists, for every \(v\), each candidate apex \(t\in N_{H_1}(v)\)
and a vertex \(w\in N_{H_1}(v)\setminus\{t\}\) not adjacent to it.
Every candidate occurs, so no neighbourhood has an apex.
There is no first fold, and \(H_1\notin\D_0\).
\end{proof}

\begin{table}[htbp]
\centering\small
\caption{Complete exclusion of an apex in every neighbourhood of \(H_1\).}
\label{4k4:tab:manual-prefix-negative}
\begin{tabular}{r l}\toprule
\(v\)&Candidate apex \(t\mapsto w\), with \(tw\notin E(H_1)\)\\\midrule
2 & $8\mapsto 13,\ 9\mapsto 10,\ 10\mapsto 9,\ 12\mapsto 9,\ 13\mapsto 8$\\
4 & $5\mapsto 8,\ 8\mapsto 5,\ 10\mapsto 12,\ 12\mapsto 10,\ 13\mapsto 8$\\
5 & $4\mapsto 9,\ 9\mapsto 4,\ 10\mapsto 9,\ 12\mapsto 9,\ 13\mapsto 12$\\
8 & $2\mapsto 4,\ 4\mapsto 2,\ 9\mapsto 4,\ 10\mapsto 9,\ 12\mapsto 9$\\
9 & $2\mapsto 5,\ 5\mapsto 2,\ 8\mapsto 5,\ 13\mapsto 8$\\
10 & $2\mapsto 4,\ 4\mapsto 2,\ 5\mapsto 2,\ 8\mapsto 5,\ 13\mapsto 8$\\
12 & $2\mapsto 4,\ 4\mapsto 2,\ 5\mapsto 2,\ 8\mapsto 5$\\
13 & $2\mapsto 4,\ 4\mapsto 2,\ 5\mapsto 2,\ 9\mapsto 4,\ 10\mapsto 9$\\
\bottomrule
\end{tabular}
\end{table}
\subsection{Certifying the prefixes and the seed}\label{4k4:app:certificates}
For the larger graphs, the neighbourhood reductions are themselves recursive. We record them as finite proof states so that every required membership and non-membership can be checked.
The main certificate contains 1,114 proof nodes for \(A\) and 184
for \(B-p\). The prefix certificates use the ordered vertex lists
\(S_1,S_2\) of this section.
All neighbourhoods and domination relations are computed in the
current induced subgraph.

\subsubsection{The recursive proof rule}
A proof state \((S,j)\), with \(j\geq-1\), asserts whether the
graph induced by \(S\) belongs to \(\D_j\).
There are five allowed forms of node.
\begin{enumerate}
\item The empty state is false at every level.
\item A non-empty cone is true at every level, with its apex supplied.
\item A non-empty non-cone is false at level \(-1\).
\item A true state with \(j\geq0\) specifies \(v\in S\), a true child
\((S\cap N(v),j-1)\), and a true child \((S\setminus\{v\},j)\).
\item A non-empty non-cone state \((S,j)\), with \(j\geq0\)
(and hence \(\card S\geq2\)), is false if it specifies, for every
\(v\in S\), a false child of at least one of those two forms.
\end{enumerate}
The checker verifies the exact subsets, levels, truth values and
child references, requires every vertex in a negative node to be
covered, and rejects cycles. Every child has fewer vertices.
It reconstructs each induced graph from the stated adjacency data,
tests every proposed apex, and checks the absence of an apex in each
negative non-cone state. A deleted vertex must belong to the current set.
The correctness of a checked proof graph therefore follows by
induction on the number of vertices, using Definition~\ref{def:hierarchy}
and Lemma~\ref{lem:inclusions} for cone states. By
the first-deletion equivalence~\eqref{eq:first}, any reduction of a state with
\(\card S\geq2\) begins with a vertex \(v\) such that
\((S\cap N(v),j-1)\) is true, and the remaining reduction shows
that \((S\setminus\{v\},j)\) is true. The false children of a
form-5 node exclude every such \(v\). This gives a finite certificate
of non-membership as well as membership, without relying on a bounded
unsuccessful search.

The three main roots are bound to
\[
(V(A),3,\mathrm{true}),\quad (V(A),2,\mathrm{false}),\quad
(V(B)\setminus\{p\},3,\mathrm{true}).
\]
The four prefix roots assert membership at levels \(1,2\) and
non-membership at levels \(0,1\), respectively.
The checker also verifies the prefix adjacencies against the
specified induced subgraphs of \(A\).
The seed and first extension are reconstructed from
Tables~\ref{4k4:tab:Aadj} and~\ref{4k4:tab:profiles}, the type graph
in~\eqref{4k4:eq:R}, and the six adjacency rules of
Section~\ref{sec:construction}.

The complete finite proof graphs are the files
\path{data/certificates/FINITE_CERTIFICATE.json.gz} and
\path{data/certificates/PREFIX_CERTIFICATES.json.gz}.
The seed and profiles are in \path{data/certificates/PARAMETERS.json};
the ordered prefix sets and their adjacencies are in
\path{data/certificates/PREFIX.json}. The checking rules are implemented in
\path{dismantling/certificates.py}. The verifier checks the supplied proof
states using integer and set operations; it does not search for reductions.

\subsubsection{The reductions and their obstructions}
For \(H_1\), delete \(a_5,a_4\) in that order; the remaining graph
is a cone with apex \(a_2\). Their current neighbourhoods belong
to \(\D_0\), as proved by the referenced child certificates.
For \(H_2\), delete \(a_6,a_5,a_4\); the current neighbourhoods
belong to \(\D_1\), and the remainder is a cone with apex \(a_2\).
For \(A=H_3\), delete \(a_0,a_1\); the current neighbourhoods belong
to \(\D_2\), and the remainder is a cone with apex \(a_{15}\).
For \(B-p\), the outer order is
\[
a_3,a_1,a_2,a_0,a_4,a_6,a_5,a_8,a_9,a_{10},a_{11}.
\]
Each current neighbourhood belongs to \(\D_2\), and the remainder
is a cone with apex \(a_7\).
These orders alone do not certify their neighbourhoods;
the complete child proof graphs do so by the preceding rules.

The negative root for \(A\notin\D_2\) has one obstruction for each
of its 16 vertices: every \(A[N_A(a_i)]\) is certified outside
\(\D_1\). The prefix negative roots similarly provide a complete
obstruction for every possible first deletion at their claimed
lower levels. A zero value of \(q\) is never used as a sufficient test of membership.

\subsection{The nineteen neighbourhood obstructions}
Table~\ref{4k4:tab:Bq} supplies all coefficients through degree seven
of every non-central neighbourhood in \(B\).
All higher degrees were computed and verified to vanish.
They are checked using the coefficient recurrence~\eqref{eq:clique-coefficients}.
The supplied checker uses this deletion recurrence to verify every
coefficient and alternating sum in the distributed certificate.
For example, the row for \(a_0\) gives
\(1-16+95-263+344-182+20=-1\).
All 19 sums are non-zero.
\begin{table}[H]
\centering\small
\begin{tabular}{r rrrrrrrr r}\toprule
\(v\)&\(c_0\)&\(c_1\)&\(c_2\)&\(c_3\)&\(c_4\)&\(c_5\)&\(c_6\)&\(c_7\)&\(q\)\\\midrule
\(a_{0}\) & 1 & 16 & 95 & 263 & 344 & 182 & 20 & 0 & $-1$\\
\(a_{1}\) & 1 & 16 & 95 & 262 & 340 & 179 & 20 & 0 & $-1$\\
\(a_{2}\) & 1 & 16 & 96 & 271 & 362 & 189 & 10 & 0 & $-7$\\
\(a_{3}\) & 1 & 16 & 95 & 261 & 335 & 172 & 16 & 0 & $-2$\\
\(a_{4}\) & 1 & 16 & 96 & 275 & 391 & 257 & 62 & 3 & $-1$\\
\(a_{5}\) & 1 & 16 & 97 & 281 & 402 & 262 & 60 & 3 & $-2$\\
\(a_{6}\) & 1 & 16 & 95 & 263 & 349 & 200 & 34 & 2 & $-2$\\
\(a_{7}\) & 1 & 16 & 97 & 279 & 388 & 232 & 39 & 1 & $-3$\\
\(a_{8}\) & 1 & 15 & 86 & 235 & 308 & 164 & 15 & 0 & $-4$\\
\(a_{9}\) & 1 & 16 & 95 & 263 & 344 & 183 & 19 & 0 & $-3$\\
\(a_{10}\) & 1 & 16 & 95 & 264 & 354 & 207 & 37 & 1 & $-1$\\
\(a_{11}\) & 1 & 14 & 75 & 194 & 249 & 143 & 27 & 2 & $-1$\\
\(a_{12}\) & 1 & 16 & 96 & 271 & 367 & 211 & 34 & 2 & $-2$\\
\(a_{13}\) & 1 & 15 & 86 & 237 & 323 & 199 & 41 & 1 & $-1$\\
\(a_{14}\) & 1 & 16 & 97 & 281 & 402 & 262 & 60 & 3 & $-2$\\
\(a_{15}\) & 1 & 16 & 96 & 272 & 375 & 230 & 48 & 3 & $-1$\\
\(x_0\) & 1 & 15 & 86 & 236 & 315 & 179 & 26 & 0 & $-2$\\
\(x_1\) & 1 & 15 & 85 & 228 & 296 & 167 & 27 & 0 & $-1$\\
\(x_2\) & 1 & 16 & 97 & 278 & 380 & 213 & 27 & 1 & $-3$\\
\bottomrule\end{tabular}
\caption{Clique coefficients of \(B[N_B(v)]\), for every non-central
vertex. All coefficients of degree greater than seven are zero.
The last column is their alternating sum.}
\label{4k4:tab:Bq}
\end{table}



\runcheck{4k4}
The output reports seven checked roots and nineteen clique rows. The root
checks certify $H_1\in\D_1\setminus\D_0$, $H_2\in\D_2\setminus\D_1$,
$A\in\D_3\setminus\D_2$, and $B-p\in\D_3$. Every coefficient in the
printed table is compared with both clique enumeration and the deletion
recurrence, including the vanishing of higher coefficients.

\subsection{The family at every level}
The certificates establish the hypotheses needed for the extension from the sixteen-vertex seed. Theorem~\ref{thm:template} then treats any number of layers.
\begin{theorem}\label{4k4:thm:family}
With the seed and profiles in this section, put $H_3=A$ and
$H_k=G_{k-3}$ for $k\geq3$, and use the two induced prefixes for $k=1,2$.
Then
\[
 H_k\in\D_k\setminus\D_{k-1},\qquad |V(H_k)|=4k+4\qquad(k\geq1).
\]
Every inclusion $H_k\subset H_{k+1}$ is induced.
\end{theorem}
\begin{proof}
The seven roots give the two prefixes and establish
$A\in\D_3\setminus\D_2$ and $B-p\in\D_3$. The nineteen non-zero sums
in Table~\ref{4k4:tab:Bq} exclude every non-central vertex as a first deletion.
Thus Theorem~\ref{thm:template} applies with $s=3$ and $d=3$.
For $k\geq3$ it gives strict level $k$ and
$16+4(k-3)=4k+4$ vertices. The prefix certificates prove the cases
$k=1,2$. The nested prefix sets and the layer rules give the induced inclusions.
\end{proof}
For $k\geq3$, the layered graphs also have
\[
 |E(G_m)|=97+16m+42m+4\binom m2+2m^2+3m(m-1)
          =97+53m+7m^2.
\]
The terms count seed edges, seed--centre edges, seed--auxiliary edges,
edges between centres or within one type, edges between the two adjacent
pairs of types, and centre--auxiliary edges, respectively.
